\documentclass{article}
\usepackage{amsmath}
\newcommand{\argmin}{\arg\!\min} 
\newcommand{\argmax}{\arg\!\max} 
\usepackage{amssymb}
\usepackage{amsfonts}
\usepackage[T1]{fontenc}
\usepackage{bm}
\usepackage{array}
\usepackage{graphicx}
\usepackage[utf8]{inputenc}
\usepackage{hyperref}
\hypersetup{
    colorlinks=true,
    linkcolor=blue,
    filecolor=magenta,
    urlcolor=cyan,
    pdftitle={Overleaf Example},
    pdfpagemode=FullScreen,
    }

\title{Cartesian space control in SMC}
\date{\today}
\author{Marko Guberina}

\begin{document}
\maketitle

\section{Introduction}
There is nothing new or special about Cartesian control in SMC.
This text is completely written in reference to
\begin{itemize}
		\item blog posts by  \href{https://scaron.info/category/robotics.html}{Stéphane Caron}, in particular:
				\begin{itemize}
						\item \href{https://scaron.info/robotics/differential-inverse-kinematics.html}{differential inverse kinematics} 
						\item \href{https://scaron.info/robotics/jacobian-of-a-kinematic-task-and-derivatives-on-manifolds.html}{Jacobian of kinematic task}
				\end{itemize}
		\item Frank C. Park, Kevin Lynch - Modern Robotics
		\item Murray, R.M., Li, Z. and Sastry, S.S., 2017. A mathematical introduction to robotic manipulation. CRC press
\end{itemize}

However, it is useful to have a summary at arms length.

\subsection{Notation used}
\begin{itemize}
		\item $ q \in \mathbb{R}^{nq}  $ - vector of joint positions
		\item $ v \in \mathbb{R}^{nv} $ - vector of joint velocities
		\item $ a\in \mathbb{R}^{nv}  $ - vector of joint acceleration
		\item $ T_{ a }^{ b } \in SE (3) $ - transformation from frame $\left\{ a \right\} $
				to $ \left\{ b \right\}   $, equivalently frame $ \left\{ b \right\}   $
				as seen in frame $ \left\{ a \right\}   $
		\item $ V \in se (3) $ - twist, i.e. 6D velocity, 
				$ V = \begin{bmatrix} v & \omega \end{bmatrix} ^{ T }  $
\end{itemize}

\section{Differential inverse kinematics}
Skipping all prerequisites on geometry and robot kinematics,
let's start with
\begin{equation}
J (q) v = V
\end{equation}
where $ J \in \mathbb{R}^{6 \times nv}  $ is the Jacobian (linearization)
of the forward kinematics map, i.e. $ J (q) = \frac{\partial FK}{\partial q}   $
and $ V  $ is end-effector twist.

Usually $ J  $ is not invertible, and we do not assume it to be invertible in this text.
Hence the task of inverse kinematics is to solve
\begin{equation}
\min_{v} \left\lVert J v - V \right \rVert ^{ 2 }
\end{equation}
We can also throw some scaling in with $ K_{ p } V  $, with
$ K_{ p } \in [0,1]  $.
If this was all there was to it, we would solve this the pseudoinverse and call it a day, i.e.
\begin{equation}
v = J^{ T } (J J^{ T } + \lambda I) V = J^{ \dagger }V
\end{equation}
But if we have, for example, linear constraints, we can interpret
the equation as a quadratic program (QP) solve it 
with a QP solver.

Moreover, we very often have additional tasks, or want to use the redundant 
joints of freedom for more graceful movement (exploit the nullspace of $ J  $).
This is especially important if there's multiple frames we care about,
for example in a mobile manipulator.
The additional tasks can be simply summed together:
\begin{equation}
\min_{v} \sum_{i}^{} w_{ i } \left\lVert J_{ i } v - K_{ i } V_{ i } \right \rVert 
\end{equation}
Solving this will result in a compromise between the two targets,
the amount of compromizing being determined with the weights $ w_{ i }  $.
Famously, this emulates the behavior of a hierarchical quadratic program (HQP).

\subsection{Inequality constraints}
The most important inequality constraint concerns joint limits,
making the problem look like
\begin{align}
		\min_{v}& \sum_{i}^{} w_{ i } 
		\left\lVert J_{ i } v - K_{ i } V_{ i } \right \rVert \\
		\text{s.t. } & v_{ \min } \leq v \leq v_{ \max }
\end{align}
The same can be done with joint angle limits 
$ q_{ \min } \leq q \leq q_{ \max }  $.

\subsection{Transcribing into standard QP form}



\subsection{Common secondary tasks}
\subsubsection{Collision avoidance via control-barrier functions}
\subsubsection{Manipulability maximization}













































\end{document}
